% !TeX root = ../main.tex
% Condensed from tables/rubric_anchors_detailed.tex; no additional scoring standard.
\begin{table*}[t]
\centering
\caption{Core anchors of the proposed rubric. Scores 1 and 5 require dimension-specific evidence; intermediate anchors and exceptions appear in Table~\ref{tab:rubric-detailed-anchors}. Calibration remains pending.}
\label{tab:rubric-core-anchors}
\begingroup\small
\setlength{\tabcolsep}{4pt}
\renewcommand{\arraystretch}{1.10}
\begin{tabular}{@{}>{\raggedright\arraybackslash}p{0.15\textwidth}>{\raggedright\arraybackslash}p{0.405\textwidth}>{\raggedright\arraybackslash}p{0.405\textwidth}@{}}
\toprule
Dimension & 1: core failure & 5: supported fulfillment \\
\midrule
D1: Specificity & Mood/style or generic dancing leaves the main movement unidentifiable. & Main movements and necessary attributes are explicit, without inventing the action. \\
D2: Body/space & Explicit body, spatial, reference-frame, or concurrent-action relations conflict. & Necessary body/spatial relations are clear and compatible, including concurrent actions. \\
D3: Local timing & Order/duration conflicts; clearly serial actions exceed available time. & Order, concurrency, sustained actions, and transitions fit the interval coherently. \\
D4: Rhythm & Claimed rhythm contradicts audible pulse/pauses, or movement rhythm is unspecified. & Repetition, pauses, speed, or movement density correspond to audible rhythm. \\
D5: Expression & Claimed force/articulation conflicts with, or depends on absent, audible features. & Force, articulation, or sustained motion responds to audible expression or grounded contrast. \\
D6: Development & Invents a key musical change, or an applicable change lacks a movement response. & Movement maintenance, emphasis, or change relates to independently identified musical development. \\
\midrule
L1: Continuity & Important boundaries impose incompatible explicit states without a resolving transition. & Important boundaries connect coherently, without decisive explicit state conflicts. \\
L2: Global structure & Unorganized concatenation or explicit conflict with audible whole-song structure. & Section functions and overall progression coherently reflect musical development. \\
L3: Motifs/variation & Unorganized repetition or unrelated changes prevent recognizable thematic development. & Recognizable motifs recur, extend, or vary at musically justified positions. \\
\bottomrule
\end{tabular}
\endgroup
\end{table*}
